% ============================================================================
% III. Методи за възстановяване
% ============================================================================
\section{Методи за възстановяване}

Методите са подредени от най-простия към най-сложния. Всеки следващ метод разчита на
по-малко информация от носителя и затова изисква повече знание за самия формат. На
практика експертът започва от първия метод и преминава към следващите само ако
предишният не е дал резултат.

\subsection{Възстановяване по метаданни на файловата система}

При изтриване записът, който сочи къде е файлът, обикновено само се отбелязва като
изтрит \cite{nist80086}. Ако този запис и данните са оцелели, файлът може да бъде
възстановен по информацията на самата файлова система. Това е
най-надеждният случай: възстановява се точно последователността от клъстери, която е
съставлявала файла, включително при фрагментиране. Ограничението е, че методът зависи
изцяло от оцелелите структури на файловата система. След форматиране или след
продължителна работа на носителя тези структури обикновено са загубени или
презаписани.

\subsection{Изрязване по сигнатури (file carving)}

\term{Изрязването на файлове} (file carving) е възстановяване на файлове само по тяхното
съдържание, без помощта на файловата система \cite{pal2009}. Програмата обхожда целия
носител и търси известни \term{сигнатури}: поредици от байтове, с които започват
(заглавие) или завършват (окончание) файловете от даден формат. Между намереното
заглавие и следващото окончание (или до определен максимален размер) байтовете се
изрязват като кандидат за файл. Инструментът Scalpel е пример за такъв подход, при който
сигнатурите се задават в конфигурационен файл, а носителят се обхожда ефективно
\cite{richard2005scalpel}.

За MP4 подходящата сигнатура е типът \code{ftyp} на байтове 4 до 7 от началото на файла,
когато \code{ftyp} е първата кутия (раздел~1.1). Окончание в същия смисъл MP4 няма. Структурата на кутиите обаче дава
по-добра възможност: от заглавието на всяка кутия се чете нейният размер, така че
програмата може да „прескача“ от кутия на кутия и да определи точния край на файла.
Това е частен случай на \term{валидиране}, разгледано в следващата точка.

Основното ограничение на изрязването по сигнатури е предположението, че файлът е записан
непрекъснато. При фрагментиран файл изрязаният участък съдържа чужди данни, а
действителното продължение на файла е другаде \cite{garfinkel2007}.

\subsection{Възстановяване на фрагментирани файлове}

Garfinkel предлага два подхода за този проблем \cite{garfinkel2007}. Първият е
\term{бързото валидиране на обекти} (fast object validation): всеки кандидат за файл се
проверява дали е формално валиден за своя формат, например дали се декодира без грешки.
Валидирането отхвърля погрешно изрязаните кандидати и позволява да се изпробват различни
комбинации от участъци. Вторият подход е \term{изрязването с откриване на празнина при
два фрагмента} (bifragment gap carving). Ако файлът е разделен на две части с чужди данни
между тях, програмата изпробва различни размери и положения на празнината, докато
валидирането не потвърди, че резултатът е цял файл.

Pal и Memon проследяват развитието на изрязването като цяло \cite{pal2009}. Общият
принцип на методите за фрагментирани файлове е един и същ: колкото по-строго може да се провери, че
даден участък е продължение на файла, толкова по-надеждно е сглобяването. За видео в
MP4 проверката може да използва две нива: дали кутиите имат съгласувани размери и дали
кодираните данни в \code{mdat} се декодират като валиден поток H.264.

\subsection{Възстановяване на ниво кадри}

Когато контейнерът е повреден или липсва, а част от данните е презаписана, търсенето на
цял файл губи смисъл. Na и съавтори предлагат друг подход: възстановяване на отделни
кадри по спецификацията на кодека \cite{na2014}. Методът разпознава кодираните кадри
директно в данните и ги декодира, като използва знанието за структурата на потока. Така
се спасява съдържанието от непрезаписаните участъци дори когато файлът като цяло не
може да бъде отворен. Както посочват авторите, методът възстановява по-голямата част от
непрезаписаното съдържание, но не и презаписаните участъци \cite{na2014}.

Този подход се опира пряко на структурата от раздел~1.2. За да се декодират кадрите, са
нужни SPS и PPS. Ако те са загубени заедно с \code{moov}, параметрите трябва да се
възстановят от друг източник, например от запис на същото устройство. Освен това
декодирането може да започне само от IDR кадър. Кадрите между началото на оцелелия
участък и първия IDR кадър в него остават недекодируеми.

\subsection{Ремонт на прекъснат запис с еталонен файл}

Прекъснатият запис от раздел~2.3 е специален случай: видеоданните са налице и
непрекъснати, липсва само \code{moov}. Инструментът \code{untrunc} решава точно тази задача
\cite{untrunc}. Той изисква \term{еталонен файл}: друг, неповреден запис от същата камера
или приложение \cite{untrunc}. Защо е нужен еталон, следва от раздел~I. На повредения
файл липсват параметрите на кодека и таблиците на пробите, а те са в \code{moov}.
Параметрите на кодека могат да се вземат от еталон, записан със същите настройки.
Размерите и отместванията на пробите трябва да се изчислят наново от самите данни в
\code{mdat}. Така по принцип се получава нова кутия \code{moov}, която прави файла
възпроизводим. Документацията на инструмента не описва вътрешната му работа и
предупреждава, че за успеха е нужен „и малко късмет“ \cite{untrunc}.

\subsection{Сравнение на методите}

Таблица~\ref{tab:methods} обобщава петте метода. Изборът между тях следва от
таблица~\ref{tab:causes}: за всяка причина за загуба има метод, който използва точно
онова, което е оцеляло.

\begin{table}[H]
\caption{Сравнение на методите за възстановяване}
\label{tab:methods}
\centering
\begin{tabular}{@{}L{3.3cm}L{3.6cm}L{3.4cm}L{3.6cm}@{}}
\toprule
\textbf{Метод} & \textbf{Какво изисква} & \textbf{Какво възстановява} & \textbf{Основно ограничение} \\
\midrule
По метаданни на файловата система & Оцелели структури на файловата система & Целия файл, вкл. фрагментиран & Не работи след форматиране \\
Изрязване по сигнатури & Само данните на носителя & Цели непрекъснати файлове & Греши при фрагментиране \\
Валидиране и сглобяване & Строга проверка на формата & Фрагментирани файлове & Изчислително скъпо \\
На ниво кадри & Параметри на кодека & Отделни кадри & Не възстановява презаписаното \\
Ремонт с еталон & Еталонен запис от същото устройство & Възпроизводим файл & Зависи от еталона \\
\bottomrule
\end{tabular}
\end{table}
